% =====================================================================
%  CHAPTER 16: 学校教学设施与设备卫生（对应大纲第十六章）
%  参考：大纲第16章
% =====================================================================
\chapter{学校教学设施与设备卫生}
\setcounter{chapter}{16}
\chapterintro{
\textbf{掌握：}校址选择与教学用房卫生要求；教室自然采光与人工照明卫生标准*；课桌椅功能尺寸及卫生管理。\newline
\textbf{熟悉：}黑板卫生；学校生活设施卫生。
}

% =====================================================================
% 第一部分：学校教学设施卫生
% =====================================================================
\section{学校教学设施卫生}

\subsection{校址选择}

\begin{keybox}
\textbf{校址选择的卫生学要求：}

\begin{enumerate}[leftmargin=*]
    \item \imp{阳光充足}：校址应选择在日照条件良好的地段，保证教学楼尤其是教室获得充足的日照。
    \item \imp{空气流通}：校址所在区域应有良好的自然通风条件，避免处于空气涡流区或静风区。
    \item \imp{地势高燥，排水通畅}：校址应选择在地势较高、地下水位较低的地段，确保排水良好，防止校园积水和潮湿。地势平坦或略有坡度，便于学生活动。
    \item \imp{交通便利}：校址应设在居民区较适中的位置，便于学生就近入学。服务半径应根据学校规模和交通条件合理确定。
    \item \imp{远离污染源和噪声源}：校址应远离交通干道（距离不少于80米）、铁路、机场、工厂、垃圾处理场等噪声和污染源。避免设在集贸市场、娱乐场所等喧闹地区附近。
\end{enumerate}

\textbf{校址选择的附加要求：}

\begin{itemize}[leftmargin=*]
    \item 校址周边环境安静，有利于教学活动的开展。
    \item 不应与不利于学生身心健康的场所（如网吧、歌舞厅等）毗邻。
    \item 应有发展余地，预留校园扩建空间。
    \item 有稳定的水源和电力供应。
\end{itemize}
\end{keybox}

\subsection{教学用房卫生要求}

\begin{keybox}
\textbf{（一）校园功能分区}

校园应按功能进行合理分区，一般分为：
\begin{itemize}[leftmargin=*]
    \item \imp{教学区}：教学楼、实验楼等。
    \item \imp{体育运动区}：操场、体育馆、游泳池等。
    \item \imp{生活区}：食堂、宿舍等。
    \item \imp{绿化区}：校园绿地、花园等。
\end{itemize}

各功能区应有明确界限，又相互联系方便，做到"动静分离"。

\textbf{（二）教学用房布局要求}

\begin{enumerate}[leftmargin=*]
    \item \imp{音乐教室、舞蹈教室}：应远离普通教学用房，避免其音响对其他教室的干扰。可有单独的声音隔离措施。
    \item \imp{南向教室}：冬至日满窗日照不小于2小时，保证教室获得充足的日照。
    \item \imp{教学楼间距}：应满足教室的日照和采光要求，一般不小于南侧建筑高度的1.5--2倍。
    \item \imp{道路系统}：校园道路应直捷、顺畅，便于人流通行和疏散。校门口应有缓冲地带，避免直接开向交通干道。
    \item \imp{升旗区}：应设有国旗升旗场地，位置醒目，便于全校师生集合。
\end{enumerate}

\textbf{（三）校园用地面积}

\begin{itemize}[leftmargin=*]
    \item \imp{中学生均用地面积}：20--25 $\mathrm{m}^2$/生。
    \item \imp{小学生均用地面积}：15--20 $\mathrm{m}^2$/生。
\end{itemize}

\textbf{（四）小学教学楼层数}

小学教学楼不宜超过4层，便于学生紧急疏散和日常活动。
\end{keybox}

\subsection{教室卫生要求}

\begin{keybox}
\textbf{教室卫生的一般要求：}

\begin{enumerate}[leftmargin=*]
    \item \imp{足够的面积}：保证每个学生有充足的活动空间，生均使用面积小学不低于1.36 $\mathrm{m}^2$，中学不低于1.39 $\mathrm{m}^2$。
    \item \imp{良好的采光照明}：保证课桌面和黑板面有足够的照度，照度分布均匀。
    \item \imp{适宜的微小气候}：教室内温度、湿度、风速适宜，空气新鲜。
    \item \imp{防止噪声干扰}：教室的隔音效果良好，外界噪声不干扰教学。
    \item \imp{便于通行和疏散}：教室门、走廊、楼梯等设计合理，便于紧急情况下迅速疏散。
    \item \imp{便于清扫}：地面和墙面材料便于清洁，有利于保持卫生。
\end{enumerate}

\textbf{教室的建筑设计要求：}

\begin{enumerate}[leftmargin=*]
    \item \imp{朝向}：教室宜朝南或南偏东/南偏西，以获得良好的日照和通风。
    \item \imp{走廊设计}：宜采用内廊或外廊设计，保证通行通畅。
    \item \imp{楼梯设计}：楼梯坡度适宜（斜度不大于30°），踏步高度一致；梯段宽度不小于疏散要求；不采用螺旋楼梯。
\end{enumerate}

\textbf{房间最小净高：}

\begin{table}[H]
\centering
\caption{各类教学用房最小净高}
\begin{tabularx}{\textwidth}{lX}
\hline
\textbf{用房类型} & \textbf{最小净高（m）} \\
\hline
小学普通教室 & 3.00 \\
初中普通教室 & 3.05 \\
高中普通教室 & 3.10 \\
舞蹈教室 & 4.50 \\
风雨操场（室内体育场） & 需满足运动安全要求（一般不低于7.00m） \\
\hline
\end{tabularx}
\end{table}
\end{keybox}

\subsection{教室的内部布置}

\begin{defbox}
\textbf{教室内部布置的卫生要求：}

\begin{enumerate}[leftmargin=*]
    \item \imp{桌椅排列}：课桌椅排列应便于学生就座和通行，前排桌缘距黑板不应小于2.0m，最后一排课桌后缘距黑板不应大于9.0m（小学）或9.5m（中学）。各列桌椅间应有不小于0.6m的纵向走道。

    \item \imp{讲台区域}：讲台长度应大于黑板长度，宽度不应小于0.8m，高度一般为0.2m，讲台两端与黑板边缘的水平距离不应小于0.2m。

    \item \imp{黑板布置}：黑板宽度小学不宜小于3.6m，中学不宜小于4.0m；黑板高度不应小于1.0m。黑板下缘距讲台面垂直距离小学宜为0.8-0.9m，中学宜为1.0-1.1m。

    \item \imp{教室环境布置}：教室内应保持整洁，墙面宜采用高亮度低彩度的装修，前墙颜色可稍暗于侧墙。教室不宜粘贴过多装饰物，以免分散学生注意力。可适当布置激励性标语和学生优秀作品展示区。
\end{enumerate}
\end{defbox}

\subsection{教室采光与照明}

\begin{keybox}
\textbf{一、自然采光}

\textbf{（一）基本卫生要求}

\begin{enumerate}[leftmargin=*]
    \item \imp{采光窗面积充足}：满足教室内各桌面获得良好的自然光照。
    \item \imp{光线方向合理}：宜双侧采光，保证光线分布均匀。单侧采光时应从学生座位左侧射入。
    \item \imp{避免眩光}：采取措施防止直射阳光产生眩光，可设置遮阳设施。
    \item \imp{室内表面反射合理}：墙壁、天花板、课桌面等的反射系数应适宜。
\end{enumerate}

\textbf{（二）自然采光评价指标}

\begin{table}[H]
\centering
\caption{教室自然采光评价指标与标准}
\begin{tabularx}{\textwidth}{l>{\raggedright\arraybackslash}p{4cm}X}
\hline
\textbf{指标} & \textbf{要求} & \textbf{说明} \\
\hline
\imp{窗地面积比} & $\geq 1:5$ & 窗户透光面积与室内地面面积之比。 \\
\imp{室深系数} & $\geq 1:2$（单侧采光） & 窗上缘高度与室深（进深）之比。 \\
\imp{入射角} & $\geq 20^\circ \sim 22^\circ$ & 室内工作面上某点与窗上缘连线与水平面所成的夹角。 \\
\imp{开角} & $\geq 4^\circ \sim 5^\circ$ & 课桌面的测定点到窗上缘连线、与到对面建筑物最高点连线的夹角差。 \\
\imp{采光系数（昼光率）} & $\geq 2\%$（课桌面） & 室内某一点的照度与同一时间室外全天空照度之比。 \\
\hline
\end{tabularx}
\end{table}

\textbf{（三）光气候系数（$K$）}

中国不同光气候分区的光气候系数：
\begin{itemize}[leftmargin=*]
    \item I区：$K = 0.85$
    \item II区：$K = 0.90$
    \item III区：$K = 1.00$
    \item IV区：$K = 1.10$
    \item V区：$K = 1.20$
\end{itemize}

光气候系数越小，表示该地区光气候条件越好。四川盆地等常年多雾地区光气候系数较高，需较大的采光口面积来满足采光要求。

\textbf{（四）室内表面反射系数}

合理的室内表面反射系数有利于提高采光效率和光线均匀度：

\begin{table}[H]
\centering
\caption{教室各表面反射系数建议值}
\begin{tabularx}{\textwidth}{lX}
\hline
\textbf{表面} & \textbf{反射系数建议值} \\
\hline
天花板 & 0.70 -- 0.80 \\
侧墙和后墙 & 0.70 -- 0.80 \\
前墙 & 0.50 -- 0.60 \\
课桌面 & 0.25 -- 0.45 \\
地面 & 0.20 -- 0.40 \\
黑板面 & 0.15 -- 0.20 \\
\hline
\end{tabularx}
\end{table}

\textbf{（五）采光均匀度}

双侧采光优于单侧采光。单侧采光时，教室远窗侧和近窗侧的照度差异较大。为减小照度不均匀度，窗间墙宽度不宜大于窗宽的1/2。

\textbf{（六）采光方向与遮阳}

\begin{itemize}[leftmargin=*]
    \item 避免阳光直射课桌面和黑板面，产生眩光。
    \item 应设置窗帘、百叶窗或遮阳板等遮阳设施。
    \item 北向教室光线稳定，但冬季日照不足；南向教室日照充足，但需注意遮阳。
\end{itemize}

\textbf{二、人工照明}

\textbf{（一）基本卫生要求}

\begin{enumerate}[leftmargin=*]
    \item 照度足够：课桌面和黑板面达到规定的照度标准。
    \item 照度均匀：教室内各点的照度差异不应过大。
    \item 避免眩光：光源不应直接或经反射面进入学生视线产生眩光。
    \item 光色适宜：光源的色温和显色性应适合教学活动需要。
    \item 安全可靠：照明设施符合电气安全标准。
\end{enumerate}

\textbf{（二）照度标准\difficulty}

\begin{table}[H]
\centering
\caption{教室人工照明照度标准}
\begin{tabularx}{\textwidth}{lX}
\hline
\textbf{指标} & \textbf{标准值} \\
\hline
课桌面照度 & $\geq 300\ \mathrm{lx}$ \\
课桌面照度均匀度 & $\geq 0.7$（最小照度/平均照度） \\
黑板面垂直照度 & $\geq 500\ \mathrm{lx}$ \\
黑板面照度均匀度 & $\geq 0.8$ \\
光源色温 & 3300--5500 K \\
\hline
\end{tabularx}
\end{table}

\textbf{（三）灯具配置要求}

\begin{itemize}[leftmargin=*]
    \item 宜采用白色荧光灯或LED照明灯具，光效高、寿命长。
    \item 灯具长轴应垂直于黑板方向布置。
    \item 灯具排列间距应均匀，保证照度均匀度。
    \item 应设置黑板灯，单独照亮黑板区域。
    \item 灯具应定期清洁和维护，及时更换老化和损坏的光源。
\end{itemize}

\textbf{（四）防眩光措施}

\begin{itemize}[leftmargin=*]
    \item 灯具应带有格栅或漫射罩，减少直接眩光。
    \item 灯具安装高度适宜，避免进入学生正常视野范围。
    \item 光源表面亮度不宜过高。
    \item 黑板面采用无光泽材料，减少反射眩光。
    \item 避免在镜面反射区布置灯具。
\end{itemize}
\end{keybox}

\subsection{教室通风与采暖}

\begin{keybox}
\textbf{（一）教室通风}

\textbf{通风方式：}

\begin{enumerate}[leftmargin=*]
    \item \imp{自然通风}：通过门窗等开口进行换气，是最基本、最经济的通风方式。应充分利用课间休息时间开窗通风。教室应设有气窗或通风小窗，便于寒冷季节进行换气。
    \item \imp{机械通风}：当自然通风不能满足要求时，应设置机械通风系统。新风量应满足卫生要求。
\end{enumerate}

\textbf{空气质量标准：}

\begin{itemize}[leftmargin=*]
    \item 教室内 $\mathrm{CO}_2$ 浓度 $\leq 0.10\%$（1000 ppm）为基本卫生要求。
    \item 教室空气中 $\mathrm{CO}_2$ 浓度是评价教室空气清洁程度的重要指标。
\end{itemize}

\textbf{最小换气次数：}

\begin{itemize}[leftmargin=*]
    \item 小学教室：不少于5次/小时。
    \item 初中教室：不少于6次/小时。
    \item 高中教室：不少于7.5次/小时。
\end{itemize}

\textbf{（二）教室采暖}

在严寒和寒冷地区，教室应设置采暖设施。

\begin{itemize}[leftmargin=*]
    \item \imp{采暖温度}：教室内温度应维持在18--22°C。
    \item \imp{温度均匀}：教室内水平温差和垂直温差不超过2°C。
    \item \imp{不产生有害物质}：采暖设施不应产生有害气体、粉尘或噪声。
    \item \imp{安全可靠}：采暖设施表面温度不宜过高，防止烫伤。
\end{itemize}
\end{keybox}

\begin{tipbox}
\textbf{通风采暖卫生标准制定的依据：}

\begin{enumerate}[leftmargin=*]
    \item \imp{CO₂浓度标准的制定依据}：教室CO₂浓度卫生标准（≤0.10\%）是基于学生脑力工作能力与空气清洁度关系的研究制定。当CO₂浓度超过0.10\%时，空气的理化性状变差（异味、细菌增殖），学生主观感觉空气污浊，注意力下降，学习效率降低。该标准综合考虑了室内外空气交换率、学生人数与教室容积比等因素。

    \item \imp{换气次数标准的制定依据}：小学不低于5次/h、初中6次/h、高中7.5次/h的换气次数标准，是基于不同学段学生的代谢率、安静时CO₂呼出量和教室容积等因素计算得出，以保证教室内CO₂浓度始终维持在0.10\%以下的卫生要求。

    \item \imp{采暖温度标准的制定依据}：教室冬季采暖温度18-22°C的标准，综合考虑了：(1)学生在静坐学习状态下的热舒适度需求；(2)不同气候区域学生的体温调节适应能力；(3)温度对学习效率和脑力工作能力的影响（过高或过低均不利于学习）；(4)中国经济发展水平和能源供应条件的可行性。
\end{enumerate}
\end{tipbox}

% =====================================================================
% 第二部分：学校教学设备卫生
% =====================================================================
\section{学校教学设备卫生}

\subsection{黑板}

\begin{keybox}
\textbf{黑板的卫生要求：}

\begin{enumerate}[leftmargin=*]
    \item \imp{表面反射系数}：黑板面应无光泽，反射系数控制在0.15--0.20，避免产生镜面反射眩光。
    \item \imp{颜色}：宜采用墨绿色或黑色，与粉笔字迹形成鲜明对比。
    \item \imp{尺寸}：黑板下缘距讲台地面高度：小学80--90 cm，中学100--110 cm。黑板宽度：小学不小于3.6 m，中学不小于4.0 m。
    \item \imp{材料}：表面应耐磨、耐擦洗，书写时不产生粉尘。鼓励使用无尘粉笔或液体粉笔。
    \item \imp{安装}：应安装在教室前墙正中央，不得有倾斜和晃动。
\end{enumerate}
\end{keybox}

\subsection{课桌椅卫生要求}

\begin{keybox}
\textbf{（一）课桌椅的基本卫生要求}

\begin{enumerate}[leftmargin=*]
    \item \imp{满足学习活动需要}：能支持书写、阅读和听课三种基本学习姿势。
    \item \imp{适合身材}：课桌椅的尺寸与学生的身高和体型相匹配，保证正确坐姿。
    \item \imp{坚固安全}：结构稳定，材料安全环保，无尖锐棱角。
    \item \imp{经济实用}：价格适中，便于维护和更换。
\end{enumerate}

\textbf{（二）课桌椅的关键功能尺寸}

\begin{table}[H]
\centering
\caption{课桌椅关键功能尺寸说明}
\begin{tabularx}{\textwidth}{lX}
\hline
\textbf{尺寸名称} & \textbf{定义与卫生意义} \\
\hline
\imp{桌高} & 桌面近胸缘距地面的高度。桌高应使学生书写时前臂能自然放在桌面上，不耸肩也不弯腰。 \\
\imp{椅高} & 椅面（座面）距地面的高度。椅高应等于或略低于小腿长，使学生足底平稳着地，大腿呈水平状态。 \\
\imp{桌椅高差} & 桌高减去椅高。高差是影响坐姿的关键尺寸。高差适宜时，学生书写时上体可自然前倾，不弯腰驼背。 \\
\imp{桌下空区} & 桌面下供放置双腿的空间。应有足够的高度和深度，使学生双腿能自由活动。 \\
\imp{桌面} & 桌面深度和宽度应满足放置书本和文具的需要。桌面宜有0°--12°的倾斜度（或水平面）。 \\
\imp{椅面} & 椅面深度（前后径）应使学生臀部得到充分支撑，但不过深压迫腘窝。 \\
\imp{椅靠背} & 靠背应在腰背部提供支撑，减轻脊柱负荷。宜采用腰靠式设计。 \\
\imp{桌椅距离} & 桌与椅之间的水平距离。负距离（椅面前缘伸入桌下）不超过4 cm为宜，可使学生在读写时保持正确姿势。\\
\hline
\end{tabularx}
\end{table}

\textbf{（三）GB/T 3976--2014 《学校课桌椅功能尺寸及技术要求》}

国家标准将课桌椅分为0--10号共11个型号，以不同颜色标识，适用身高范围从 $\leq 119\ \mathrm{cm}$ 到 $\geq 187.5\ \mathrm{cm}$。

\begin{table}[H]
\centering
\caption{GB/T 3976--2014 课桌椅型号（简化）}
\begin{tabularx}{\textwidth}{llXl}
\hline
\textbf{型号} & \textbf{适用身高范围（cm）} & \textbf{颜色标识} & \textbf{桌面高（mm）} \\
\hline
0号  & $\geq 187.5$ & 紫色 & 810 \\
1号  & 180.0 -- 187.4 & 蓝色 & 780 \\
2号  & 172.5 -- 179.9 & 绿色 & 750 \\
3号  & 165.0 -- 172.4 & 橙色 & 720 \\
4号  & 157.5 -- 164.9 & 黄色 & 690 \\
5号  & 150.0 -- 157.4 & 红色 & 660 \\
6号  & 142.5 -- 149.9 & 白色 & 630 \\
7号  & 135.0 -- 142.4 & 紫色 & 600 \\
8号  & 127.5 -- 134.9 & 蓝色 & 570 \\
9号  & 120.0 -- 127.4 & 绿色 & 540 \\
10号 & $\leq 119.0$ & 橙色 & 510 \\
\hline
\end{tabularx}
\end{table}

\textbf{（四）课桌椅的卫生管理}

\begin{enumerate}[leftmargin=*]
    \item \imp{合理配置}：每间教室至少配置2种以上型号的课桌椅，满足不同身高学生的需要。
    \item \imp{定期调整}：每学期根据学生身高变化调整课桌椅分配。
    \item \imp{合格率要求}：课桌椅与学生身高的符合率达到80\%以上。
    \item \imp{科学排座}：身高较低者坐在前排，身高较高者坐在后排，并定期轮换。
\end{enumerate}
\end{keybox}

% =====================================================================
% 第三部分：学校生活设施卫生
% =====================================================================
\section{学校生活设施卫生}

\subsection{学校食堂}

\begin{keybox}
\textbf{（一）食堂选址}

\begin{itemize}[leftmargin=*]
    \item 应远离污染源（厕所、垃圾站、有毒有害场所等），距离不少于25 m。
    \item 应有良好的给排水条件。
    \item 便于食品运输和学生就餐。
\end{itemize}

\textbf{（二）食堂功能分区}

学校食堂一般分为三大功能区：

\begin{enumerate}[leftmargin=*]
    \item \imp{食品加工区}：包括原料处理间、切配间、烹调间、面点间、洗消间等。应生熟分开，防止交叉污染。
    \item \imp{就餐区}：餐厅面积应满足学生就餐需要。中小学食堂就餐座位数一般不少于学生人数的1/3（分批就餐）。
    \item \imp{辅助区}：包括更衣室、库房、办公室、卫生间等。
\end{enumerate}

\textbf{（三）食堂卫生管理}

\begin{itemize}[leftmargin=*]
    \item 严格执行食品卫生法律法规。
    \item 从业人员持有健康合格证。
    \item 落实食品留样制度（每份留样不少于125g，留样48小时）。
    \item 餐具严格消毒。
    \item 定期消毒和灭蝇灭鼠。
\end{itemize}
\end{keybox}

\subsection{校车}

\begin{defbox}
学校应配备符合安全标准的校车，保障学生上下学交通安全。
\end{defbox}

\subsection{其他生活设施}

\begin{keybox}
\textbf{（一）供水设施}

\begin{itemize}[leftmargin=*]
    \item 学校应为学生提供充足、安全、方便的饮用水。
    \item 饮用水水质符合国家生活饮用水卫生标准（GB 5749）。
    \item 供水设备定期清洗消毒，水质定期检测。
    \item 直饮水设备需取得卫生许可批件，有定期更换滤芯和消毒记录。
    \item 不提倡长期饮用纯净水（缺乏矿物质）。
\end{itemize}

\textbf{（二）厕所卫生}

\begin{itemize}[leftmargin=*]
    \item 厕所数量应满足学生使用需要。教学楼每层均应设厕所。
    \item 女生厕所蹲位数应多于男生。
    \item 厕所应通风良好，采光充足，便于清洁。
    \item 应配备洗手设施和洗手用品（肥皂或洗手液）。
    \item 定期清扫和消毒，保持清洁无异味。
    \item 化粪池应远离水源，有防渗漏措施。
\end{itemize}

\textbf{（三）学生宿舍}

\begin{itemize}[leftmargin=*]
    \item 宿舍朝向宜朝南或偏南，保证日照和通风。
    \item 生均居住面积不低于3 $\mathrm{m}^2$。
    \item 室内净高不低于3.00 m（寒冷地区不低于2.80 m）。
    \item 应设置独立的卫生间或公共盥洗室和淋浴间。
    \item 宿舍应有专人管理，建立健全的安全管理制度。
    \item 宿舍应定期消毒，保持清洁卫生。
    \item 保证学生充足的睡眠时间，制定合理的作息制度。
    \item 定期对学生进行宿舍安全教育（防火、防盗、防传染病等）。
\end{itemize}

\textbf{（四）其他公共设施}

\begin{itemize}[leftmargin=*]
    \item 学校应设有医务室（保健室），配备基本医疗设备和常用药品。
    \item 学校应有心理健康辅导室，面积不少于30 $\mathrm{m}^2$。
    \item 学校应有足够的垃圾收集设施，实行垃圾分类。
    \item 校园绿化覆盖率不低于30\%。
    \item 校园内的道路和活动场地应有足够的夜间照明。
\end{itemize}
\end{keybox}

\newpage
